\documentclass[12pt]{article}
\usepackage[english]{babel}
\usepackage[english]{isodate}
\usepackage[table,svgnames]{xcolor}
\usepackage{cmbright}
\usepackage{url}
\usepackage[utf8x]{inputenc}
\usepackage[T1]{fontenc}
\usepackage{longtable}
\usepackage{amsmath}
\usepackage{graphicx}
\usepackage{parskip}
\usepackage{fancyhdr}
\usepackage{vmargin}
\usepackage{hyperref}
\usepackage{pgfgantt}
\usepackage{xparse}
\usepackage{float}
\usepackage{tabularx}
\usepackage{titling}
\usepackage[nogroupskip,toc,acronym]{glossaries}
\usepackage{fancyhdr}
\usepackage[authoryear]{natbib}
\usepackage{enumitem}
%\bibliographystyle{fontysIEEtranN}

\pagestyle{fancy}
\fancyhf{}
\lhead{\includegraphics[width=0.2\linewidth]{TreewatchLogo.pdf}}
\rhead{Personal dossier René Karoff}
\rfoot{Page \thepage}

% alternate row colors for all tables
\definecolor{lightGrey}{rgb}{0.95,0.95,0.95}

\let\oldtable\table
\let\endoldtable\endtable
\renewenvironment{table}{\rowcolors{2}{lightGrey}{}\oldtable}{\endoldtable}

\let\oldtabular\tabular
\let\endoldtabular\endtabular
\renewenvironment{tabular}{\rowcolors{2}{lightGrey}{}\oldtabular}{\endoldtabular}

\let\oldtabularx\tabularx
\let\endoldtabularx\endtabularx
\renewenvironment{tabularx}{\rowcolors{2}{lightGrey}{}\oldtabularx}{\endoldtabularx}

\let\oldlongtable\longtable
\let\endoldlongtable\endlongtable
\renewenvironment{longtable}{\rowcolors{2}{lightGrey}{}\oldlongtable} {\endoldlongtable}

\graphicspath{{../img/}}

\renewcommand{\arraystretch}{1.5}

\DeclareDocumentCommand{\newdualentry}{ O{} O{} m m m m } {
  \newglossaryentry{gls-#3}{name={#5},text={#5\glsadd{#3}},
    description={#6},#1
  }
  \makeglossaries
  \newacronym[see={[Glossary:]{gls-#3}},#2]{#3}{#4}{#5\glsadd{gls-#3}}
}


\newlist{SMART}{description}{2}
\setlist[SMART]{leftmargin=12em,style=nextline}

\newlist{STARR}{description}{2}
\setlist[STARR]{leftmargin=8em,style=nextline}


%\loadglsentries{../glossary.tex}

\setmarginsrb{3 cm}{2.5 cm}{3 cm}{2.5 cm}{1 cm}{1.5 cm}{1 cm}{1.5 cm}

%\makeglossaries
\begin{document}

%%%%%%%%%%%%%%%%%%%%%%%%%%%%%%%%%%%%%%%%%%%%%%%%%%%%%%%%%%%%%%%%%%%%%%%%%%%%%%%%%%%%%%%%% Preface of the report

    \pagenumbering{roman} % Roman numerals for page counter


    % Title Page
    \begin{titlingpage}
        \begin{center}
            \begin{minipage}{\linewidth}
            \centering
            %University logo
            \includegraphics[width=0.3\linewidth]{FontysLogo.pdf}
            \par
            \vspace{3cm}
            %Thesis title
            {\uppercase
                {\Large Personal dossier\\ René Karoff \\ 2015 Sofa GTL \\ TreeWatch Project
            \par
            \vspace{3cm}}}
            \includegraphics[width=0.5\linewidth]{TreewatchLogo.pdf}
            \par
            \vspace{2cm}
            %Author's name
            {René Karoff\\
	         \par}
            \vspace{2cm}

            %Date
            \today
            \end{minipage}
        \end{center}
    \end{titlingpage}
    \clearpage

\setmarginsrb{3 cm}{2.5 cm}{3 cm}{2.5 cm}{1.5 cm}{1.5 cm}{1 cm}{1.5 cm}
    \section*{Document information}
\addcontentsline{toc}{section}{Document information}
\renewenvironment{tabular}{\oldtabular}{\endoldtabular}
	\begin{tabular}{ll}
		\textbf{Document name:} & Personal dossier René Karoff\\
		\textbf{Document owner:} & René Karoff \\
		\textbf{Company/Organisation:} & Fleuren Baarlo \\
		\textbf{Contact person:} & René Karoff \\
		\textbf{Date:} & \today \\
		\textbf{Place:} & Fontys University of Applied Science Venlo \\
		\textbf{Author:} & \parbox[t]{5cm}{
		René Karoff\\ r.karoff@student.fontys.nl\\ 2198664 \\}
	\end{tabular}
\renewenvironment{tabular}{\rowcolors{2}{lightGrey}{}\oldtabular}{\endoldtabular}

    %\printglossary[type=\acronymtype]
    %\printglossary
    \pagebreak

    %\listoffigures
    %\addcontentsline{toc}{section}{\listfigurename}
    %\listoftables
    %\addcontentsline{toc}{section}{\listtablename}
    %\pagebreak

    \tableofcontents
    \clearpage

%%%%%%%%%%%%%%%%%%%%%%%%%%%%%%%%%%%%%%%%%%%%%%%%%%%%%%%%%%%%%%%%%%%%%%%%%%%%%%%%%%%%%%%%% Main Part of the Report

    \pagenumbering{arabic}
    
    \section{Improve quality management skills}
	\subsection{SMART learning goal}
	\begin{SMART}
	    \item[Specific] As a project quality manager in the GreenTechLab Software Factory, my goal is to increase my ability to identify and specify quality standards as well as to assure the quality and to control it.
	    \item[Measurable] Since I had issues with the quality management in my internship report, I will receive feedback from the other team members during the project using a questionnaire, as well as from the evaluation board at the end of this project. Furthermore I could get in touch with the customer to see if he was satisfied with the quality management.
	    \item[Attainable] Since we spoke about Quality Management in the second year course Business Management, I have several documents to rely on and to work from. Furthermore there is a bunch of documentation about this topic available on the internet.
	    \item[Relevant] In my future career, in which I want to become a project manager. As a project manager it is also important to know the quality management.
	    \item[Time-limited]By January 2016, this project will be closed, therefore I will be able to manage the quality of a project until then.
	    \item[My complete goal] To improve my quality management by January 2016.
	\end{SMART}
	
	%This requirement was drafted when we thought we were going to develop using web technology and not via xamarin. So the requirement has changed from web based development to normal software development. \\
	%I choose this one, because as a software engineer it is always good to keep improving your software skills, especially in new environments. For me it was new to develop for multi platform, so there is a lot to learn.
	
	\subsection{STARR Analysis}
	\begin{STARR}
	    \item[Situation] At the start of the project I was elected as the project quality manager.
	    \item[Task] My task was to make sure that everyone sticked to the quality standard which were set in the beginning of the project.
	    \item[Action] In the beginning of the project I started to set up several quality guidelines regarding documentation and coding conventions.
	    \item[Result] The result of my role as the quality manager of this project was a quality management plan (REFERENCE-QMPLAN) in which the standards were defined. Furthermore it was stated how the group should be able to reach the targeted quality standards. One example is that a documentation template was created, so that every document which was created had the same overall layout, as well as titlepage and a document information page. (REFERENCE-TEMPLATE).
	    \item[Reflection] I wanted to be quality manager because of the issues with the quality management in my internship report. Since there was a lack in quality management I decided to go for this role. The project started well, the standards were defined fast and everyone agreed with the quality management plan. But then during the project i recognized that there were issues with one of the defined standards. I told the group in each situation that we should write tests but the testing was only poorly done. Now, 
	\end{STARR}
	
	\section{Improve coding skills}
	This
	\begin{SMART}
		\item[Specific]
		\item[Measurable]
		\item[Attainable]
		\item[Relevant]
		\item[Time-limited]
		\item[My complete goal]
	\end{SMART}
	
	\begin{STARR}
		\item[Situation]
		\item[Task]
		\item[Action]
		\item[Result]
		\item[Reflection]
	\end{STARR}
	
	\section{Improve communication skills}
	\subsection{SMART learning goal}
	\begin{SMART}
	    \item[Specific] As a project quality manager in the GreenTechLab Software Factory, my goal is to increase my ability to speak to a customer more freely in english.
	    \item[Measurable] I will receive feedback from the team members using a questionnaire. Furthermore I will receive feedback from our lecturer.
	    \item[Attainable] We will have several meetings with our customer, and as a project quality manager I will talk to the customer about some of the requirements and about some quality standards.
	    \item[Relevant] In my future career, in which I want to become a project manager, speaking freely to customers is a core competence.
	    \item[Time-limited] By January 2016, this project will be closed, therefore I’ll have to higher my communication skills until then
	    \item[My complete goal] To improve my communication skills by January 2016.
	\end{SMART}
	
	\subsection{STARR Analysis}
	\begin{STARR}
	    \item[Situation]
	    \item[Task]
	    \item[Action]
	    \item[Result]
	    \item[Reflection]
	\end{STARR}
	
	\section{Improve analytical skills}
	
	\subsection{SMART learning goal}
	\begin{SMART}
	    \item[Specific] As a project quality manager in the GreenTechLab Software Factory, my goal is to improve my analytical competencies in new situations. Meaning that everytime a new situation occurs, it has to analyzed.
	    \item[Measurable] I will receive feedback from the team members using a questionnaire.
	    \item[Attainable] Since we are working within a scrum environment, it is possible that the requirements change completely or at least the prioritization, therefore I would have to react and to analyze the new situation.
	    \item[Relevant] In my future career, in which I want to become a project manager, good analytical competencies are necessary to be successful.
	    \item[Time-limited] By January 2016, this project will be closed, therefore I’ll have to higher my analytical competencies by then.
	    \item[My complete goal] To improve my analytical competencies by January 2016.
	\end{SMART}
	
	\subsection{STARR Analysis}
	\begin{STARR}
	    \item[Situation]
	    \item[Task]
	    \item[Action]
	    \item[Result]
	    \item[Reflection]
	\end{STARR}
	
	
	%\clearpage
	
	%\bibliography{biblist}
	
	\end{document}
